系统受理一份购买，便承担了一段跨进程、跨保存域的责任。入口的 HTTP 并发只描述链路开始处的资源，C/T、prepared 库存、EV 和 dispatch 意图分别占用后续阶段的空间与处理预算。容量设计沿这条事实链分开计算流入、等待和完成，再决定何处限制新工作、何处为恢复保留资源。

\subsection{为每种责任建立独立计量}

设入口每秒收到 5000 次调用，其中 2000 次首次建立 $S$，3000 次读取已有操作。付款平均每份需要 1 次发送和 0.2 次查询，正常平台尝试速率便约为 $2000\times1.2=2400$ 次每秒。若发布确认丢失使平均每份 EV 发布 1.05 次，消息写入速率是每秒 2100 条。消费事务按 EV 与业务守卫汇合重复消息，实际首次履约速率最多接近每秒 2000 份。这些数字是同一条链的不同计量，分别适用于入口、平台、broker 和业务效果。

\CTwentyOneFigure{ch21-05}{业务责任和通信尝试分别消耗阶段资源，积压在持久保存位置累积}{fig:ch21-loads}

未决 T 表示尚未结束的付款执行责任，未发布 EV 积压表示已经形成、尚未完成发布确认的事件。发布以后尚无持久消费交接证据的 EV 另列一项。平台恢复以后，大批 T 结束，同时产生 EV，积压从付款阶段移到传播阶段。消费者一侧又分为 broker 尚未交付、已交付未确认、消费事务排队以及 dispatch 未知。每一层分别记录数量、字节、最老年龄、流入和流出，值守者据此把内存、磁盘和外部额度关联到实际等待的责任。

库存还要按件和商品计量。一万笔 prepared 事务若都占用同一热门 SKU，与占用一万种冷门 SKU 的影响很不相同。退款用金额约束恢复：原收款 3000 分，成功退款 1000 分，尚未知的退款占额 1500 分，再次受理的最大金额为 500 分。未知退款转到人工处理以后，原 RF 继续保存这 1500 分占额，财务取得确定结果后再由权威调整占用。

在稳定观察窗口内，设某阶段到达速率为 $\lambda$，平均停留时间为 $W$，平均在途数量为 $N$，可用 Little 定律估算：
\[
N=\lambda W\text{。}
\]
每秒 100 份付款、平均 2 秒取得确定结果，对应平均约 200 份在途责任。若其中 1\% 因平台异步处理平均等待 600 秒，其余 99\% 等待 1 秒，则平均停留时间为 $0.99\times1+0.01\times600=6.99$ 秒，在途责任约 699 份。少数长尾已经占据大部分存量，按正常 1 秒延迟规划 100 个工作槽便会低估责任规模。

等待中的责任可以落盘，不必各占一个线程。工作者只在到期时领取有限批次，发出查询以后释放执行槽，下一次动作时间保存在责任或持久调度结构中。内存批次按数量和字节双重受限，系统启动时逐页扫描到期索引。这样，数十万份等待责任主要占持久空间，实际并发由平台、连接池和回填能力控制。

\subsection{正常流与恢复流共用哪几类资源}

一份首次购买会经历受理提交、资格领取、结果回填和发布标记等权威命令，批量方式与实现协议决定实际次数。设本例平均每份需要 4 次组内命令，每秒 2000 份新购买便产生约 8000 次命令。已有身份查询每秒 3000 次，其中 80\% 从不可覆盖的原结果缓存或入口合并观察得到回答，剩余每秒 600 次访问权威。权威的设计输入因此为每秒约 8000 次命令和 600 次读取，库存事务与恢复命令再分别加入预算。

平台能力以调用和并发两种限制给出。若平台允许每秒 3000 次调用，平均占用连接 120 毫秒，正常每秒 2400 次尝试对应平均 288 个在途调用。选择 400 个并发槽可以留出波动余量。若响应缩短到 50 毫秒，400 个槽每秒可轮换约 8000 次，发送器就会撞上每秒调用额度。速率令牌按时间补充，每次尝试消耗一份，令牌积累上限约束允许的短时突发；已经发出而长期未返的调用则继续占用连接。发送器同时核验速率、在途数和每个商户范围的配额，分别控制单位时间压力和长尾存量。

一批未知付款的查询通常还要访问权威、解析平台响应并回填。若平台每秒能回答 1000 次，权威回填每秒最多完成 600 份，恢复完成率就受后者限制。若调度器分配每秒 600 次查询，每份未知责任平均先经过 1.5 次查询才确定，则这份查询预算对应每秒约 400 份完成。预算用完成的责任数计算排空，用尝试次数计算平台压力，两者通过查询成功比例连接。

磁盘也要沿阶段算。设每份新增责任及索引的逻辑写入为 2 KiB，业务结果与事件再增加 3 KiB，每秒 2000 份对应约 9.77 MiB 的逻辑新增数据。三副本复制、日志、快照与索引会增加实际写入和占用。若消息暂停 10 分钟，120 万份事件每份 1.5 KiB，逻辑载荷约 1.72 GiB，还要加消息索引和副本。规划按磁盘可用余量及写放大预留空间，恢复期间同时下载快照、整理表和回放日志会争用同一磁盘。

恢复流还可能占满控制和身份服务。地域切换后，入口同时刷新目录，工作者同时重新领取责任，客户端同时重连。可以让刷新请求在进程内合并，使用有界并发和带抖动的重试，按分区分批恢复后台扫描。此时控制组面对的是配置恢复流，数据组面对的是执行恢复流，两者分别设置池和速率，避免把一次缓存失效放大成所有组件同步重建。

\subsection{热点倾斜决定分区的实际余量}

负载分布足够均匀时，可以先用总流入除以分区数估算单分区压力。设十个分区，每个健康时能处理每秒 1000 个本例命令，总流入每秒 6000 个。若 b1 承担 40\%，它收到每秒 2400 个，另外九个组平均只有 400 个。集群总能力尚有余量，b1 的队列却每秒增长约 1400 条。此时扩容要落到 b1 的承载资源、流入限制或业务分区方式，空闲分区的能力才能通过重新分配工作被使用。

热点来自不同层次。同一 $S$ 的重试可以合并读取，相同账户的不同 $S$ 则可能竞争同一余额与约束，热门 SKU 的库存预留又形成跨账户共享点。扩大执行线程会让更多请求争用同一状态。对第一类热点，合并与结果缓存能减少重复工作；对第二类，按账户或资源排队、限制在途新责任和合理分区更有效；对第三类，可按业务允许的库存分配结构，由权威保存互不重叠的地域额度，各地在获配额度内受理。分配者把已用数量、未用额度与尚未分配的库存一同计量，回收未用额度则经过下文的委托结算。

以 b1 的同身份突发为例，100 个连接同时询问 K1，入口进行一次权威观察，其余等待该观察的结果。若 100 次分别购买不同商品，则形成 100 个独立 $S$，权威为各自成功受理的身份保存独立绑定与实际 P，每份都参与余额、库存与准入。账户维度用于安排公平份额，完整身份用于合并同一操作的观察，资源关联和操作类型则用于定位余额或库存热点。

在线迁移也有容量前提。把 b1 搬到另一组，解决的是当前承载组的资源分布，b1 内部的串行资源竞争继续存在。若源和目标都在恢复排空，完整图像复制还会占用网络和磁盘。迁移调度器根据目标余量、prepared 责任和图像大小决定启动时机。迁移完成后，入口刷新与缓存重热又形成短期突发，放量速度纳入同一预算。

\subsection{重试怎样把等待变成额外负载}

客户端每隔 1 秒重试一次，一个 30 秒才结束的操作可能产生约 30 次观察。每秒 100 份新操作持续到达、平均都等 30 秒时，稳态约有 3000 份在途操作，每秒轮询量约为 3000 次。若每次轮询又触发一次平台查询，原来每秒 100 份业务就放大成约 3000 次平台尝试。入口把进度读取和平台恢复调度分开，客户端查询读取已保存的最近观察，后台按到期时间查询平台。

多层自动重试还会相乘。外层 HTTP 库最多 3 次、业务 SDK 最多 3 次、工作者每轮最多 3 次，若层间缺少统一期限，某次故障可能触发 27 次下游发送。稳定原键可以把这些发送汇合到同一业务效果，连接、带宽和平台额度却已经被消耗。调用链选择明确的重试负责人，其他层只对尚未发送的连接建立错误作有界处理，并传播剩余期限。

退避给系统争取恢复空间。设初次等待上限为 $d_0$，最大等待上限为 $d_{\max}$，重试序号 $k$ 从 0 开始。第 $k$ 次等待上限为
\[
d_k=\min(d_{\max},d_0\,2^k)\text{，}
\]
实际等待在 $[0,d_k]$ 内均匀随机取值，避免所有工作者在 1、2、4 秒整点重新聚集。进程重启后从持久的下一次动作时间继续，重试次数和期限用于调度、告警与转人工。一个付款达到 20 次尝试时仍可能未知，调度器把它转到人工核对队列，并保留 C、原键和占额。

服务器可在忙响应中给出建议等待时间。调用方保持原 $S$ 与原输入，避免每次退避都生成新身份。对已发送但结果未知的受理命令，入口返回未知并提供原身份查询路径。只有明确在创建责任之前拒绝的请求，才能被描述为本次未受理。忙、未知和业务拒绝分别对应不同事实，客户端据此选择原身份恢复或修改业务意图。

\subsection{准入落在建立持久责任的边界}

购买提交入口先完成调用者认证和授权，再核验会话存在、处于活动状态，以及可信时间尚未达到原固定截止。任一资格条件不成立，旧 POST 返回 \code{operation\_expired}。通过这道公共入口以后，权威才查询完整 $S$：若 $A_J$ 已存在，核对完整输入，恢复原结果或返回进度；若尚未受理，再检查新购买的业务约束和新增责任额度，并在同一仲裁中创建 $A_J$、C/T 和相应配额占用。图~\ref{fig:ch21-admission} 的两条分支分别计量已有操作和新责任。会话关闭以后，授权查询入口和后台 C/T 恢复继续读取已受理事实，使用各自的恢复预算。

\CTwentyOneFigure{ch21-06}{提交先核验公共资格，通过后已有身份恢复原绑定，新身份原子消耗责任额度}{fig:ch21-admission}

设 b1 最多允许 1000 份未决付款，当前为 999。两个入口各自读取 999，再各创建一份，会得到 1001。可由分区状态机在同一条受理命令中检查并增加计数，结果结束时原子释放占用。多维准入还可以检查未发布字节、prepared 件数、退款金额和平台范围。计数是权威事实的派生索引，恢复时可与责任明细核对，避免计数漂移使系统永久误拒或放宽上限。

另一种方式是向入口委托互不重叠的额度。它减少逐请求协调，却增加委托记录、已用额度结算与失联处理。例如入口得到 10 份额度，过期前已经受理 6 份，完整使用历史可以让控制者保留 6 份责任占用、回收另外 4 份。若入口失联时，控制者只确认了 6 份已用、2 份未用，另外 2 份尚未回传，便先回收明确未用的 2 份，将剩余 2 份列为待结算，取得使用历史以后再处理。跨地域委托还把时钟、配置传播及停止使用旧额度的时刻纳入协议。

背压沿积压方向传递。消息域暂停使 EV 字节逼近阈值时，先限制可能继续增加 EV 的新购买。已经付款的订单由发布恢复排空，原结果查询继续。平台故障使 T 达到上限时，新责任减少，原键查询和确定结果回填优先获得预算。准入响应速度越快，越能避免新连接和请求体在入口堆积。

请求体在读取过程中就受尺寸限制。等待队列同时限制条数与字节，100 个含 1 MiB 附件的请求已占约 100 MiB，另有解析对象、连接缓冲和追踪状态。大附件可先存不可变对象，在规范化业务输入中引用内容摘要与版本；它的上传、权限和保留另有协议。业务去重比较读取保存的语义输入，日志摘要仅用于定位。

\subsection{合并等待和公平恢复调度}

通过提交入口资格检查以后，合并表按完整 $S$ 建立，并在加入等待前比较输入。第一个观察者读取权威，其余连接共享已观察结果。表内同时保存等待者数、创建时间和截止时间，每个 $S$ 及整个进程都有上限。超过上限时，入口返回忙或结果未知，并给出按完整 $S$ 授权查询的位置。入口已经从权威取得绑定时，可以同时附上实际 P/C/T；权威读取尚未返回或受理回执丢失时，响应保留未知，不填写候选编号。客户端保留原输入、E1 和 K，稍后从同一身份继续。入口重启丢失合并表以后，重新访问权威便能恢复原操作。

\CTwentyOneFigure{ch21-07}{内存中的合并观察减少重复工作，连接退出以后原业务责任继续保存}{fig:ch21-coalesce}

合并覆盖同一入口内同时等待的观察。若请求分散到十个地域入口，最多可能有十次权威读取。共享缓存或边缘路由可以继续减少读取，代价是增加一次网络访问与缓存失效管理。某个等待者取消 HTTP 时，入口释放它的连接和订阅，其他等待者继续接收结果。若用户同时申请撤销，权威读取原 C：尚未发送的责任按撤销规则关闭，已经产生未知效果的责任先查询平台，已确定收款的购买则沿退款路径处理。

恢复调度采用有限并发池，并把新发送、未知查询、结果回填、TX 决定恢复和事件发布分成可计量类别。可为每类保留最低速率和并发，将暂时空闲份额借给其他类别。平台恢复时，旧未知付款不会因新购买源源不断而长期排在后面。数据库回填池与平台查询池也分开，避免一批慢 HTTP 占满所有能提交确定结果的执行线程。

\CTwentyOneFigure{ch21-08}{恢复按类别和分区分配资源，单个热点受有限批次与并发约束}{fig:ch21-fairness}

简单轮转可以每轮从 b1、b2、b3 各取有限份到期责任。若不同责任代价悬殊，按条数平均会失衡。可以给分区积累调度额度，查询一次扣除估算的调用成本，大对象发布扣除字节成本，长调用又受独立在途上限限制。调度器按最早未决时间、下一次动作时间和证据期限选取本分区工作。即将失去平台查询证据的 C 得到更高优先级，同时为普通 EV 保留最低份额。

公平性需要可观察的结果。每个分区记录等待到领取的延迟、已借与保留预算、到期未处理数量和最老年龄。b1 每秒产生大量工作时，b2 的少量未知付款也应在其最低份额内得到处理。将所有未决对象按创建时间全表排序，一旦最早的一批持续超时，就可能让后面的分区长期得不到调度。

\subsection{用稳定条件计算故障后的排空}

设当前积压为 $B$ 份，恢复完成率为 $\mu_{\mathrm{rec}}$，同一队列继续流入的速率为 $\lambda_{\mathrm{res}}$，两者在计算窗口内近似稳定。净排空率为 $\mu_{\mathrm{rec}}-\lambda_{\mathrm{res}}$。只有它为正时，排空时间才能近似写为
\[
t_{\mathrm{drain}}=\frac{B}{\mu_{\mathrm{rec}}-\lambda_{\mathrm{res}}}\text{。}
\]
积压 6000 份，每秒完成 150 份，同时进入 50 份，约需 60 秒。若每秒完成 50 份且又进入 50 份，恢复流达到平衡，存量保持不降。降低新责任流入、提高该阶段完成能力或改变工作成本，才能获得排空余量。

再算一条完整链，假定各阶段能力在窗口内稳定，付款恢复每份都能终结，且每份终结产生一份 EV。平台暂停 10 分钟，期间每秒受理 200 份，积压为 12 万份 T。恢复后保留正常新流每秒 200 份，平台能每秒终结 500 份，则 T 约在 $120000/(500-200)=400$ 秒后排空。付款终结每秒产生约 500 份 EV。发布能力为每秒 600 条，可以跟上此时的事件流；消费能力若为每秒 350 份，400 秒内便新增长约 6 万份消费积压。T 排空以后，消费以每秒 $350-200=150$ 份的净速度继续恢复，再经过约 400 秒排空。整条链的恢复终点约在第 800 秒。

\CTwentyOneFigure{ch21-19}{故障积压从付款移向消费，端到端恢复由各阶段净排空能力共同决定}{fig:ch21-drain}

实际调度从区间计数估计完成率，付款拒绝、重复查询和事件比例变化分别进入流量换算，队列斜率与最老年龄用于调整准入和恢复份额。节点失效时，幸存副本同时承担领导者转移、快照和缓存重热。容量规划据单域失效后的可用能力分配这些工作：先预留原结果读取、确定结果回填和决定恢复的份额，再由剩余能力确定新购买流入。

容量观测把新 $S$ 比例、同身份重试、热点集中度、平台长尾、消息暂停时长和消费者处理成本关联到对应阶段。目录刷新变慢而提交正常时，值守者检查缓存命中、并发刷新和控制组压力。单组提交变慢时，沿热点、复制磁盘和在途命令定位。T 增长而权威正常时，比较平台尝试率与责任完成率；T 下降且 EV 增长时，比较发布完成率与消息域余量；发布确认正常而 I/B/F 停滞时，再检查消费事务、业务守卫冲突和 dispatch 调度。各阶段的队列变化给出下一步调整预算的位置。
